\documentclass[spanish,aspectratio=169]{beamer}

% ==================== PAQUETES ESENCIALES ====================
\usepackage[utf8]{inputenc}
\usepackage[spanish,es-tabla]{babel}
\usepackage[T1]{fontenc}
\usepackage{csquotes}
\usepackage{graphicx}
\usepackage{booktabs}
\usepackage{array}
\usepackage{url}
\usepackage{hyperref}
\usepackage{xcolor}
\usepackage{tikz}
\usepackage{tcolorbox}

% ==================== CONFIGURACIÓN DEL TEMA ====================
\usetheme{Madrid}
\usecolortheme{default}

% Colores personalizados (sobrios y profesionales)
\definecolor{azulpucp}{RGB}{0,51,102}
\definecolor{grisoscuro}{RGB}{64,64,64}
\definecolor{azulclaro}{RGB}{230,240,250}

\setbeamercolor{palette primary}{bg=azulpucp,fg=white}
\setbeamercolor{palette secondary}{bg=azulpucp,fg=white}
\setbeamercolor{palette tertiary}{bg=azulpucp,fg=white}
\setbeamercolor{palette quaternary}{bg=azulpucp,fg=white}
\setbeamercolor{structure}{fg=azulpucp}
\setbeamercolor{section in toc}{fg=azulpucp}
\setbeamercolor{block title}{bg=azulpucp,fg=white}
\setbeamercolor{block body}{bg=azulclaro,fg=black}

% ==================== INFORMACIÓN DEL DOCUMENTO ====================
\title[Estilo APA 7ma Ed.]{Sistema APA de Citación y Referencias}
\subtitle{7ma Edición - Guía Práctica para Estudiantes}
\author{Centrum PUCP}
\institute{Centro de Negocios - Pontificia Universidad Católica del Perú}
\date{\today}

% ==================== CONFIGURACIÓN DE ENCABEZADO/PIE ====================
\setbeamertemplate{navigation symbols}{}
\setbeamertemplate{footline}[frame number]

% ==================== INICIO DEL DOCUMENTO ====================
\begin{document}
	
	% ==================== PORTADA ====================
	\begin{frame}
		\titlepage
	\end{frame}
	
	% ==================== TABLA DE CONTENIDOS ====================
	\begin{frame}{Contenido}
		\tableofcontents
	\end{frame}
	
	% ==================== SECCIÓN 1: ESTRUCTURA ====================
	\section{Estructura del Documento Académico}
	
	\begin{frame}{Estructura General según APA}
		\begin{block}{Componentes Principales}
			Un documento académico en estilo APA debe contener:
		\end{block}
		
		\begin{enumerate}
			\item \textbf{Portada} (Título, autor, afiliación institucional)
			\item \textbf{Resumen/Abstract} (150-250 palabras)
			\item \textbf{Cuerpo del documento}
			\begin{itemize}
				\item Introducción
				\item Método
				\item Resultados
				\item Discusión
			\end{itemize}
			\item \textbf{Referencias}
			\item \textbf{Notas al pie de página} o notas al final (opcional)
			\item \textbf{Tablas}
			\item \textbf{Figuras}
			\item \textbf{Apéndices}
		\end{enumerate}
	\end{frame}
	
	\begin{frame}{Portada}
		\begin{block}{Elementos de la Portada}
			\begin{itemize}
				\item Título del trabajo (máximo 12 palabras)
				\item Nombre completo del autor(es)
				\item Afiliación institucional
				\item Curso/Programa
				\item Nombre del docente
				\item Fecha de presentación
			\end{itemize}
		\end{block}
		
		\vspace{0.3cm}
		\textbf{Nota importante:} En trabajos de investigación final y tesis, la portada se enumera pero \alert{no se muestra el número}. La siguiente página inicia con numeración romana en minúscula (ii).
	\end{frame}
	
	\begin{frame}{Resumen Ejecutivo/Abstract}
		\begin{block}{Características del Resumen}
			\begin{itemize}
				\item Extensión: 150-250 palabras
				\item Un solo párrafo
				\item Sin sangría en la primera línea
				\item Debe incluir: objetivo, método, resultados principales y conclusiones
				\item Palabras clave: 3-5 términos
			\end{itemize}
		\end{block}
		
		\begin{alertblock}{Diferencia Resumen/Abstract}
			\begin{itemize}
				\item \textbf{Resumen Ejecutivo:} en español
				\item \textbf{Abstract:} en inglés (traducción del resumen)
			\end{itemize}
		\end{alertblock}
	\end{frame}
	
	% ==================== SECCIÓN 2: FORMATO ====================
	\section{Formato del Documento}
	
	\begin{frame}{Formato General del Documento}
		\begin{columns}[T]
			\begin{column}{0.48\textwidth}
				\begin{block}{Configuración de Página}
					\begin{itemize}
						\item \textbf{Tamaño de papel:} A4
						\item \textbf{Márgenes:} 1" o 2.54 cm (superior, inferior, derecho, izquierdo)
						\item \textbf{Numeración:} lado superior derecho
					\end{itemize}
				\end{block}
			\end{column}
			
			\begin{column}{0.48\textwidth}
				\begin{block}{Tipografía}
					\begin{itemize}
						\item \textbf{Fuente recomendada:} Times New Roman 12 pts.
						\item \textbf{Interlineado:} doble en todo el documento
						\item \textbf{Sangría:} 1.27 cm o 0.5"
					\end{itemize}
				\end{block}
			\end{column}
		\end{columns}
		
		\vspace{0.5cm}
		\begin{exampleblock}{Excepciones al formato general}
			\begin{itemize}
				\item Pie de página: interlineado sencillo
				\item Tablas/figuras: interlineado simple, 1.5 o doble (según mejor presentación)
				\item Figuras: fuente sin serifas (Calibri, Arial) entre 8-14 pts.
			\end{itemize}
		\end{exampleblock}
	\end{frame}
	
	\begin{frame}{Encabezados y Niveles}
		\begin{block}{Sistema de 5 Niveles de Encabezados}
			\small
			\begin{enumerate}
				\item \textbf{Nivel 1:} Centrado, Negritas, Mayúsculas y Minúsculas
				\item \textbf{Nivel 2:} Justificado a la Izquierda, Negritas, Mayúsculas y Minúsculas
				\item \textbf{Nivel 3:} Justificado a la Izquierda, \textit{Cursivas}, Negritas, Mayúsculas y Minúsculas
				\item \textbf{Nivel 4:} Sangría, Negritas, Mayúsculas y Minúsculas con Punto Seguido.
				\item \textbf{Nivel 5:} Sangría, Negritas, \textit{Cursivas}, Mayúsculas y Minúsculas con Punto Seguido.
			\end{enumerate}
		\end{block}
		
		\begin{alertblock}{Recomendaciones}
			\begin{itemize}
				\item Numerar títulos y subtítulos en trabajos finales y tesis
				\item Evitar ``subtítulos viudos'' (al final de página sin contenido)
				\item No tener solo un subtítulo bajo un nivel (dos o más, o ninguno)
			\end{itemize}
		\end{alertblock}
	\end{frame}
	
	\begin{frame}{Seriación y Viñetas}
		\begin{columns}[T]
			\begin{column}{0.48\textwidth}
				\begin{block}{Dentro de un Párrafo}
					Los elementos en serie se identifican con letras en minúsculas entre paréntesis:
					
					\vspace{0.2cm}
					Los pasos son: (a) identificar el problema, (b) analizar las causas, y (c) proponer soluciones.
				\end{block}
				
				\begin{block}{Numeraciones}
					Se usan cuando hay orden o importancia:
					
					\vspace{0.2cm}
					1. Primer paso\\
					2. Segundo paso\\
					3. Tercer paso
				\end{block}
			\end{column}
			
			\begin{column}{0.48\textwidth}
				\begin{block}{Viñetas}
					Para elementos sin orden específico:
					
					\vspace{0.2cm}
					\begin{itemize}
						\item Elemento uno
						\item Elemento dos
						\item Elemento tres
					\end{itemize}
				\end{block}
				
				\begin{alertblock}{Importante}
					\begin{itemize}
						\item Usar el mismo estilo de viñeta en todo el documento
						\item Iniciar al nivel de sangría: 0.5"
					\end{itemize}
				\end{alertblock}
			\end{column}
		\end{columns}
	\end{frame}
	
	% ==================== SECCIÓN 3: ESTILO ====================
	\section{Estilo y Convenciones}
	
	\begin{frame}{Uso de Mayúsculas y Minúsculas}
		\begin{block}{Reglas Generales}
			\begin{itemize}
				\item \textbf{Títulos de encabezados:} Primera letra de cada palabra principal en mayúscula
				\item \textbf{Títulos en referencias:} Solo primera letra del título y subtítulo
				\item \textbf{Nombres propios:} Siempre con mayúscula
				\item \textbf{Meses:} en minúscula en español, mayúscula en inglés
			\end{itemize}
		\end{block}
		
		\begin{exampleblock}{Ejemplos}
			\begin{itemize}
				\item \textbf{Encabezado:} Metodología de la Investigación
				\item \textbf{Referencia libro:} El proceso estratégico: Un enfoque de gerencia
				\item \textbf{Fecha español:} 8 de diciembre de 2021
				\item \textbf{Fecha inglés:} May 13, 2021
			\end{itemize}
		\end{exampleblock}
	\end{frame}
	
	\begin{frame}{Uso de Cursivas}
		\begin{block}{Se usa cursiva para:}
			\begin{itemize}
				\item Títulos de libros, publicaciones periódicas, películas, programas de TV
				\item Títulos de tablas y figuras
				\item Términos clave cuando se introducen por primera vez
				\item Palabras en otros idiomas (excepto términos latinos comunes)
			\end{itemize}
		\end{block}
		
		\begin{exampleblock}{Ejemplos}
			\begin{itemize}
				\item \textit{Tabla 1}
				\item \textit{Figura 2}
				\item El concepto de \textit{empowerment} es fundamental...
				\item El término \textit{zeitgeist} se refiere a...
				\item Pero: et al., a priori, vis-à-vis (NO en cursiva)
			\end{itemize}
		\end{exampleblock}
	\end{frame}
	
	\begin{frame}{Uso de Comillas}
		\begin{block}{Se usan comillas para:}
			\begin{itemize}
				\item Citas textuales de menos de 40 palabras
				\item Títulos de artículos o capítulos
				\item Ejemplos lingüísticos
				\item Comentarios irónicos o expresiones inventadas
			\end{itemize}
		\end{block}
		
		\begin{exampleblock}{Ejemplo de Cita Textual Corta}
			\small
			El líder servicial ``demuestra que está al pendiente de su equipo, es empático con cada integrante, los ayuda a desarrollarse y los motiva según sus intereses'' (Rojas, 2025, p. 5).
		\end{exampleblock}
		
		\begin{alertblock}{Importante}
			NO usar comillas en citas de 40 o más palabras (usar bloque sangrado)
		\end{alertblock}
	\end{frame}
	
	\begin{frame}{Números}
		\begin{block}{Reglas para Números}
			\begin{itemize}
				\item \textbf{Números del 1 al 9:} en letras (excepto en mediciones, estadísticas, etc.)
				\item \textbf{Números 10 en adelante:} en cifras
				\item \textbf{Decimales:} usar punto (.)
				\item \textbf{Miles:} usar coma (,)
				\item \textbf{Millones:} usar apóstrofo (')
			\end{itemize}
		\end{block}
		
		\begin{exampleblock}{Ejemplos}
			\begin{itemize}
				\item cinco participantes, 12 estudiantes
				\item 3.14 (decimal)
				\item 1,500 (mil quinientos)
				\item 2'500,000 (dos millones quinientos mil)
				\item El 75\% de la muestra...
			\end{itemize}
		\end{exampleblock}
	\end{frame}
	
	\begin{frame}{URLs y DOIs}
		\begin{block}{Digital Object Identifier (DOI)}
			\begin{itemize}
				\item Es un código alfanumérico único asignado por la editorial
				\item Preferir DOI sobre URL cuando esté disponible
				\item Formato: https://doi.org/10.xxxx/xxxxx
				\item No usar ``Recuperado de'' ni fecha de recuperación (salvo excepciones)
			\end{itemize}
		\end{block}
		
		\begin{exampleblock}{Ejemplo en Referencias}
			\small
			Yelin, J. R. (2017). Breve estado de la cuestión animal. \textit{Perífrasis. Revista de Literatura, Teoría y Crítica, 8}(15), 29-43. https://doi.org/10.25025/perifrasis201781502
		\end{exampleblock}
		
		\begin{alertblock}{URL}
			Los cortes automáticos de URL por el procesador de texto son válidos
		\end{alertblock}
	\end{frame}
	
	\begin{frame}{Abreviaturas Comunes}
		\begin{columns}[T]
			\begin{column}{0.48\textwidth}
				\begin{block}{Abreviaturas Latinas}
					\begin{itemize}
						\item \textbf{et al.} = y otros (3+ autores)
						\item \textbf{p.} = página
						\item \textbf{pp.} = páginas
						\item \textbf{párr.} = párrafo
						\item \textbf{ed.} = edición
						\item \textbf{Vol.} = Volumen
						\item \textbf{No.} = Número
					\end{itemize}
				\end{block}
			\end{column}
			
			\begin{column}{0.48\textwidth}
				\begin{block}{Otras Abreviaturas}
					\begin{itemize}
						\item \textbf{s.f.} = sin fecha
						\item \textbf{comp.} = compilador
						\item \textbf{coord.} = coordinador
						\item \textbf{ed. acad.} = editor académico
						\item \textbf{trad.} = traductor
					\end{itemize}
				\end{block}
				
				\begin{alertblock}{Importante}
					et al., a priori, etc. NO llevan cursiva (son términos latinos incorporados al español académico)
				\end{alertblock}
			\end{column}
		\end{columns}
	\end{frame}
	
	% ==================== SECCIÓN 4: CITAS ====================
	\section{Citas en el Texto}
	
	\begin{frame}{Tipos de Citas}
		\begin{block}{Dos Tipos Principales}
			\begin{enumerate}
				\item \textbf{Citas Textuales (Directas):} Reproducen exactamente las palabras del autor
				\begin{itemize}
					\item Menos de 40 palabras: entre comillas
					\item 40 o más palabras: en bloque sangrado
				\end{itemize}
				\item \textbf{Citas Parafraseadas (Indirectas):} Redacción propia de las ideas del autor
		\end{enumerate}
	\end{block}
	
	\begin{alertblock}{Recomendación Académica}
		\begin{itemize}
			\item Preferir paráfrasis sobre citas textuales
			\item Las citas textuales solo cuando sea esencial mantener la idea original exacta
			\item Todas las citas requieren su referencia completa al final
		\end{itemize}
	\end{alertblock}
\end{frame}

\begin{frame}{Estilos de Citación: Parentética vs. Narrativa}
	\begin{block}{Citación Parentética (basada en el texto)}
		El autor aparece \alert{dentro} del paréntesis:
		
		\vspace{0.2cm}
		\small
		Los sistemas de aguas residuales constituyen una de las infraestructuras urbanas más importantes (Hernández Rodríguez \& Torres, 2019).
	\end{block}
	
	\begin{block}{Citación Narrativa (basada en el autor)}
		El autor aparece \alert{fuera} del paréntesis:
		
		\vspace{0.2cm}
		\small
		Hernández Rodríguez y Torres (2019) indicaron que los sistemas de aguas residuales constituyen una de las infraestructuras urbanas más importantes.
	\end{block}
	
	\begin{alertblock}{Nota}
		Ambos estilos son correctos. Usar según lo que se desee enfatizar.
	\end{alertblock}
\end{frame}

\begin{frame}{Formato de Citas según Número de Autores}
	\small
	\begin{block}{Tabla de Formatos}
	\begin{tabular}{p{3.5cm} l l}  % Cambia a p{} para la col1 para permitir multiline
		\toprule
		\textbf{Autores} & \textbf{Citación Parentética} & \textbf{Citación Narrativa} \\
		\midrule
		Un autor & (Zurita, 2025) & Zurita (2025) \\
		\midrule
		Dos autores & (Pizarro \& Valdivia, 2024) & Pizarro y Valdivia (2024) \\
		\midrule
		Tres o más & (Arque et al., 2023) & Arque et al. (2023) \\
		\midrule
		Grupo con abreviatura (1ra vez) & (Instituto Nacional de Estadística e Informática [INEI], 2024) & Instituto Nacional de Estadística e Informática (INEI, 2024) \\
		\midrule
		Grupo (después) & (INEI, 2024) & INEI (2024) \\
		\midrule
		Grupo sin abreviatura & (Centrum PUCP, 2025) & Centrum PUCP (2025) \\
		\bottomrule
	\end{tabular}
	\end{block}
\end{frame}

\begin{frame}{Citas Textuales: Menos de 40 Palabras}
	\begin{block}{Características}
		\begin{itemize}
			\item Entre comillas
			\item Incorporada en el párrafo
			\item Incluir número de página
			\item Punto final \alert{después} del paréntesis
		\end{itemize}
	\end{block}
	
	\begin{exampleblock}{Ejemplo Real del Manual}
		\small
		El líder servicial ``demuestra que está al pendiente de su equipo, es empático con cada integrante, los ayuda a desarrollarse y los motiva según sus intereses'' (Rojas, 2025, p. 5).
	\end{exampleblock}
	
	\begin{alertblock}{Variante con Autor en Narrativa}
		\small
		Según Rojas (2025), el líder servicial ``demuestra que está al pendiente de su equipo, es empático con cada integrante'' (p. 5).
	\end{alertblock}
\end{frame}

\begin{frame}[fragile]{Citas Textuales: 40 o Más Palabras}
	\begin{block}{Características}
		\begin{itemize}
			\item Sin comillas
			\item Bloque sangrado a 0.5" (1.27 cm)
			\item Punto \alert{antes} del paréntesis de citación
			\item Doble espacio mantenido
		\end{itemize}
	\end{block}
	
	\begin{exampleblock}{Ejemplo del Manual (formato aproximado)}
		\footnotesize
		Esta es la sangría del párrafo previo. La citación en bloque empieza al nivel de la sangría:
		
		\hspace{0.5cm}\begin{minipage}{0.85\textwidth}
			Un líder busca servir, por ello demuestra que está al pendiente de su equipo, es empático con cada integrante, los ayuda a desarrollarse y los motiva según sus intereses. Además, conoce muy bien a cada uno de ellos, no porque los desea liderar sino porque los valora como personas. (Rojas, 2025, p. 5)
		\end{minipage}
	\end{exampleblock}
\end{frame}

\begin{frame}{Citas Parafraseadas (Indirectas)}
	\begin{block}{Características}
		\begin{itemize}
			\item Redacción con palabras propias
			\item Mantener la idea del autor
			\item Dar crédito a la fuente
			\item No requiere comillas
			\item Número de página opcional (recomendado si se quiere dirigir al lector)
		\end{itemize}
	\end{block}
	
	\begin{exampleblock}{Ejemplo}
		\small
		Los líderes serviciales se caracterizan por estar atentos a las necesidades de su equipo, mostrando empatía y ayudando en su desarrollo personal y profesional (Rojas, 2025).
	\end{exampleblock}
	
	\begin{alertblock}{Recomendación}
		Este tipo de cita es la más valorada académicamente porque demuestra análisis y pensamiento crítico
	\end{alertblock}
\end{frame}

\begin{frame}{Citas Secundarias}
	\begin{block}{Definición}
		Cuando se cita a un autor cuya obra original no se pudo consultar, pero se encontró citado en otro autor (fuente secundaria).
	\end{block}
	
	\begin{exampleblock}{Formato}
		\small
		Valdez (2025, como se citó en Sánchez, 2023) propuso que...
	\end{exampleblock}
	
	\begin{alertblock}{Importante}
		\begin{itemize}
			\item En Referencias solo se incluye la fuente secundaria (Sánchez, 2023)
			\item Evitar este tipo de citas siempre que sea posible
			\item Preferir consultar la fuente original
		\end{itemize}
	\end{alertblock}
\end{frame}

\begin{frame}{Comunicaciones Personales}
	\begin{block}{¿Qué son?}
		Entrevistas personales, correos electrónicos, WhatsApp, mensajes de texto, llamadas telefónicas, conferencias no grabadas, etc.
	\end{block}
	
	\begin{exampleblock}{Formato de Citación}
		\small
		\begin{itemize}
			\item Citación narrativa: N. Apellido (comunicación personal, 26 de junio de 2025)
			\item Citación parentética: (N. Apellido, comunicación personal, 26 de junio de 2025)
		\end{itemize}
	\end{exampleblock}
	
	\begin{alertblock}{Importante}
		\begin{itemize}
			\item Se citan en el texto pero \alert{NO} se incluyen en Referencias
			\item No son recuperables por el lector
			\item Usar solo cuando sea necesario
		\end{itemize}
	\end{alertblock}
\end{frame}

\begin{frame}{Recomendaciones Adicionales para Citas}
	\begin{block}{Buenas Prácticas}
		\begin{itemize}
			\item Citar siempre al autor original (evitar citas secundarias)
			\item Usar tiempo pasado para presentar ideas del autor
			\item No usar coma entre sujeto y verbo: \alert{Incorrecto:} Valdez (2019), indicó
			\item No sobrecargar con citas del mismo autor en cada oración
			\item Evitar tener muy pocas citas (riesgo de plagio)
			\item Escribir correctamente nombres y apellidos de autores
		\end{itemize}
	\end{block}
	
	\begin{alertblock}{Plagio y Autoplagio}
		\begin{itemize}
			\item No citar = Plagio (voluntario o involuntario)
			\item No citar trabajos propios previos = Autoplagio
		\end{itemize}
	\end{alertblock}
\end{frame}

% ==================== SECCIÓN 5: REFERENCIAS ====================
\section{Referencias}

\begin{frame}{Lista de Referencias: Aspectos Generales}
	\begin{block}{Características}
		\begin{itemize}
			\item Título: ``Referencias'' (NO Bibliografía)
			\item Orden alfabético por apellido del primer autor
			\item Sangría francesa de 1.27 cm (0.5")
			\item Doble espacio entre referencias
			\item Solo incluir fuentes citadas en el texto
		\end{itemize}
	\end{block}
	
	\begin{alertblock}{Regla Fundamental}
		\Large
		\centering
		Toda cita debe tener su referencia\\
		Toda referencia debe tener su cita
	\end{alertblock}
\end{frame}

\begin{frame}{Formato General de Referencias}
	\begin{block}{Estructura Básica}
		\small
		Autor, A. A. (año). \textit{Título del documento}. Editorial o página web.
	\end{block}
	
	\begin{block}{Elementos Clave}
		\begin{itemize}
			\item \textbf{Autor:} Apellido, Inicial(es) del nombre
			\item \textbf{Año:} Entre paréntesis
			\item \textbf{Título:} En cursiva para obras principales (libros, revistas)
			\item \textbf{Título:} Entre comillas para partes (artículos, capítulos)
			\item \textbf{Editorial:} Sin términos como ``Editorial'', ``Inc.'', ``S.A.''
			\item \textbf{DOI o URL:} Al final, sin punto después
		\end{itemize}
	\end{block}
	
	\begin{exampleblock}{Autores Múltiples}
		\begin{itemize}
			\item Hasta 20 autores: listar todos separados por coma
			\item 21 o más: primeros 19 + ... + autor final
		\end{itemize}
	\end{exampleblock}
\end{frame}

\begin{frame}{Referencias: Libro Impreso}
	\begin{block}{Formato}
		\small
		Apellido, N. M. (año). \textit{Título del libro: Subtítulo}. Editorial.
	\end{block}
	
	\begin{exampleblock}{Ejemplo - Un Autor}
		\footnotesize
		Senge, P. (1992). \textit{La quinta disciplina}. Granica.
	\end{exampleblock}
	
	\begin{exampleblock}{Ejemplo - Dos Autores}
		\footnotesize
		Londoño, R., \& Restrepo, G. (1995). \textit{Diez historias de vida: ``Las Marías''}. Fundación Social.
	\end{exampleblock}
	
	\begin{exampleblock}{Ejemplo - Tres o Más Autores}
		\footnotesize
		O'Byrne Orozco, M. C., Quintana Guerrero, Í., \& Daza Caicedo, R. (2018). \textit{La obra arquitectónica de Le Corbusier}. Ediciones Uniandes.
	\end{exampleblock}
\end{frame}

\begin{frame}{Referencias: Libro Digital}
	\begin{block}{Formato}
		\small
		Apellido, N. M. (año). \textit{Título del libro: Subtítulo}. Editorial. URL o DOI
	\end{block}
	
	\begin{exampleblock}{Ejemplo del Manual}
		\footnotesize
		Lamprea Montealegre, E. (2019). \textit{El derecho de la naturaleza: Una aproximación interdisciplinaria a los estudios ambientales}. Ediciones Uniandes. https://uniandes.ipublishcentral.com/product/el-derecho-de-la-naturaleza
	\end{exampleblock}
	
	\begin{alertblock}{Nota}
		\begin{itemize}
			\item No incluir fecha de recuperación (salvo excepciones)
			\item No usar ``Recuperado de'' antes de URL
			\item Preferir DOI sobre URL cuando esté disponible
		\end{itemize}
	\end{alertblock}
\end{frame}

\begin{frame}{Referencias: Capítulo de Libro}
	\begin{block}{Formato}
		\small
		Apellido, N. M. (año). Título del capítulo. En N. M. Apellido (Ed.), \textit{Título del libro} (pp. xx-xx). Editorial.
	\end{block}
	
	\begin{exampleblock}{Ejemplo del Manual de Chicago}
		\footnotesize
		Fernández, F. (2005). La guerra naval después de la era vikinga. En M. Keen (Ed.), \textit{Historia de la guerra en la Edad Media} (pp. 295-321). Océano.
	\end{exampleblock}
	
	\begin{block}{Nota sobre pp.}
		\begin{itemize}
			\item Incluir rango completo de páginas del capítulo
			\item No abreviar: pp. 295-321 (NO pp. 295-21)
		\end{itemize}
	\end{block}
\end{frame}

\begin{frame}{Referencias: Artículo de Revista Académica}
	\begin{block}{Formato}
		\small
		Apellido, N. M. (año). Título del artículo. \textit{Nombre de la Revista, Volumen}(Número), pp-pp. DOI o URL
	\end{block}
	
	\begin{exampleblock}{Ejemplo del Manual - Con DOI}
		\footnotesize
		Yelin, J. R. (2017). Breve estado de la cuestión animal. \textit{Perífrasis. Revista de Literatura, Teoría y Crítica, 8}(15), 29-43. https://doi.org/10.25025/perifrasis201781502
	\end{exampleblock}
	
	\begin{exampleblock}{Ejemplo del Manual - Sin DOI}
		\footnotesize
		Mazzitelli, C., \& Garrido Domené, F. (2019). Las variedades del español a través del doblaje cinematográfico. \textit{Anuario de Letras, 7}(2), 63-82. https://doi.org/10.19130/iifl.adel.7.2.2019.1549
	\end{exampleblock}
\end{frame}

\begin{frame}{Referencias: Artículo de Periódico}
	\begin{block}{Formato}
		\small
		Apellido, N. M. (año, día de mes). Título del artículo. \textit{Nombre del Periódico}. URL
	\end{block}
	
	\begin{exampleblock}{Ejemplo - Con Autor}
		\footnotesize
		Román, V., \& Castilhos, W. (2019, 18 de noviembre). Los recortes en ciencia sacuden a América Latina. \textit{El País}. https://elpais.com/elpais/2019/11/15/ciencia/1573816247\_796813.html
	\end{exampleblock}
	
	\begin{exampleblock}{Ejemplo - Sin Autor}
		\footnotesize
		Aprueban política de transformación digital e inteligencia artificial en el país. (2019, 13 de noviembre). \textit{El Espectador}. https://www.elespectador.com/tecnologia/aprueban-politica...
	\end{exampleblock}
\end{frame}

\begin{frame}{Referencias: Tesis}
	\begin{block}{Formato}
		\small
		Apellido, N. M. (año). \textit{Título de la tesis} [Tesis de maestría/doctorado, Institución]. Repositorio o URL
	\end{block}
	
	\begin{exampleblock}{Ejemplo del Manual}
		\footnotesize
		Eljaiek Avendaño, H. J. (2015). \textit{Cuerpo y sujeto político en el ballet: Tensiones en el proceso formativo} [Tesis de maestría, Universidad Distrital]. http://hdl.handle.net/11349/6271
	\end{exampleblock}
	
	\begin{exampleblock}{Ejemplo - Tesis Impresa}
		\footnotesize
		Bermúdez, M. E. (2009). \textit{Sinergia y aprendizaje: Un modelo escolar de contacto entre grupos y aprendizaje a través del servicio} [Tesis de maestría inédita]. Universidad de los Andes.
	\end{exampleblock}
\end{frame}

\begin{frame}{Referencias: Documentos Legales}
	\begin{block}{Formato General - Ley}
		\small
		Nombre de la ley, Número. (Fecha). Título completo. \textit{Diario Oficial} N° xxxx.
	\end{block}
	
	\begin{exampleblock}{Ejemplo - Ley}
		\footnotesize
		Ley 1450 de 2011. (16 de junio). Por la cual se expide el Plan Nacional de Desarrollo 2010-2014. \textit{Diario Oficial} 48.102.
	\end{exampleblock}
	
	\begin{exampleblock}{Ejemplo - Decreto}
		\footnotesize
		Decreto 1072 de 2015. (26 de mayo). Por medio del cual se expide el Decreto Único Reglamentario del Sector Trabajo. Ministerio del Trabajo. https://cutt.ly/HeL5AIf
	\end{exampleblock}
\end{frame}

\begin{frame}{Referencias: Página Web}
	\begin{block}{Formato}
		\small
		Apellido, N. M. (año, día de mes). Título de la página. \textit{Nombre del Sitio Web}. URL
	\end{block}
	
	\begin{exampleblock}{Ejemplo del Manual}
		\footnotesize
		Guzmán Hormazábal, L. (2019, 15 de noviembre). Urbanización cambia los microbios en las casas y en los humanos. \textit{SciDevNet}. https://cutt.ly/1eL4rEZ
	\end{exampleblock}
	
	\begin{exampleblock}{Ejemplo - Organización}
		\footnotesize
		Organización Panamericana de la Salud. (2019, 21 de octubre). \textit{25 años sin polio en las Américas}. https://www.paho.org/col/index.php?option=com\_content\&view=article\&id=3273
	\end{exampleblock}
\end{frame}

\begin{frame}{Referencias: Redes Sociales y Multimedia}
	\begin{block}{Formato - Twitter/X}
		\small
		Autor [@usuario]. (año, día de mes). Contenido del tweet [Tweet]. Twitter. URL
	\end{block}
	
	\begin{exampleblock}{Ejemplo}
		\footnotesize
		Movilidad Bogotá [@SectorMovilidad]. (2019, 19 de octubre). Te esperamos en la avenida Suba con cra. 80 (21 Ángeles) a partir de las 1:00 p.m. para que realices tu \#RegistroBici [Tweet]. Twitter. https://twitter.com/SectorMovilidad/status/1196799287591608320
	\end{exampleblock}
	
	\begin{exampleblock}{Video de YouTube}
		\footnotesize
		Aula 365. (2016, 1 de diciembre). \textit{Qué es la contaminación ambiental} [Video]. YouTube. https://www.youtube.com/watch?v=TV-YEQOIFuQ
	\end{exampleblock}
\end{frame}

% ==================== SECCIÓN 6: TABLAS Y FIGURAS ====================
\section{Tablas y Figuras}

\begin{frame}{Tablas y Figuras: Conceptos Básicos}
	\begin{block}{Definiciones}
		\begin{itemize}
			\item \textbf{Tablas:} Presentación de datos en filas y columnas
			\item \textbf{Figuras:} Todo lo demás (gráficos, diagramas, fotografías, mapas, ilustraciones, infografías, etc.)
		\end{itemize}
	\end{block}
	
	\begin{alertblock}{Importante}
		\begin{itemize}
			\item Solo existen ``Tablas'' y ``Figuras''
			\item NO usar términos como: Cuadro, Gráfico, Foto, Diagrama, etc.
			\item Cada tabla/figura debe tener su llamado en el texto
			\item Numeración correlativa: Tabla 1, Tabla 2, Figura 1, Figura 2
		\end{itemize}
	\end{alertblock}
	
	\begin{block}{Listas Requeridas}
		\begin{itemize}
			\item Lista de Tablas
			\item Lista de Figuras
		\end{itemize}
	\end{block}
\end{frame}

\begin{frame}{Estructura de Tablas y Figuras}
	\begin{block}{Componentes (de arriba hacia abajo)}
		\begin{enumerate}
			\item \textbf{Número:} Tabla 1 / Figura 1 (en negritas, sin cursiva)
			\item \textbf{Título:} En cursiva, mayúsculas y minúsculas, alineado a la izquierda
			\item \textbf{Contenido:} La tabla o figura propiamente dicha
			\item \textbf{Nota} (opcional): Para aclaraciones
			\item \textbf{Fuente/Citación:} Información completa de la fuente
		\end{enumerate}
	\end{block}
	
	\begin{alertblock}{Formato de Fuente}
		\begin{itemize}
			\item \textbf{Adaptado de:} cuando se modificó la tabla/figura original
			\item \textbf{Tomado de:} cuando se presenta sin modificaciones
		\end{itemize}
	\end{alertblock}
\end{frame}

\begin{frame}[fragile]{Ejemplo de Tabla}
	\begin{exampleblock}{Modelo Completo}
		\footnotesize
		\textbf{Tabla 1}
		
		\textit{Características Demográficas de los Participantes}
		
		\begin{center}
			\begin{tabular}{lcc}
				\toprule
				Característica & Frecuencia & Porcentaje \\
				\midrule
				Género & & \\
				\quad Masculino & 45 & 52.3 \\
				\quad Femenino & 41 & 47.7 \\
				Edad promedio & 34.5 años & -- \\
				\bottomrule
			\end{tabular}
		\end{center}
		
		\textit{Nota.} Los datos corresponden a la muestra total (N = 86).
		
		\textit{Nota.} Adaptado de ``Título del Artículo,'' por N. M. Apellido, año, \textit{Nombre de la Revista, Volumen}(Número), p. XX. Copyright YYYY por Editorial.
	\end{exampleblock}
\end{frame}

\begin{frame}{Ejemplo de Figura}
	\begin{exampleblock}{Modelo de Figura}
		\footnotesize
		\textbf{Figura 1}
		
		\textit{Proceso de Toma de Decisiones Estratégicas}
		
		\vspace{0.3cm}
		\begin{center}
			\begin{tikzpicture}[scale=0.7, every node/.style={transform shape}]
				\node[draw, rectangle, minimum width=2cm, minimum height=0.8cm] (a) at (0,0) {Análisis};
				\node[draw, rectangle, minimum width=2cm, minimum height=0.8cm] (b) at (3,0) {Decisión};
				\node[draw, rectangle, minimum width=2cm, minimum height=0.8cm] (c) at (6,0) {Implementación};
				\draw[->, thick] (a) -- (b);
				\draw[->, thick] (b) -- (c);
			\end{tikzpicture}
		\end{center}
		
		\vspace{0.2cm}
		\textit{Nota.} Este diagrama ilustra las tres fases principales del proceso.
		
		\textit{Nota.} Tomado de \textit{Título del Libro}, por N. M. Apellido, año, p. XX. Editorial. Copyright YYYY por Autor/Editorial.
	\end{exampleblock}
\end{frame}

\begin{frame}{Recomendaciones para Tablas y Figuras}
	\begin{block}{Buenas Prácticas}
		\begin{itemize}
			\item Usar fuente más pequeña dentro de tablas/figuras si es necesario (mínimo 8 pts.)
			\item Para figuras: usar fuente sin serifas (Calibri, Arial, Lucida Sans Unicode)
			\item Evitar dejar espacios en blanco antes/después
			\item No cortar tablas entre páginas (usar página horizontal si es necesario)
			\item Presentar en página nueva si no cabe en la actual
			\item Mantener el llamado cerca del contenido
		\end{itemize}
	\end{block}
	
	\begin{alertblock}{Citación Obligatoria}
		Toda tabla o figura tomada o adaptada de una fuente externa debe citarse completamente. La citación debe aparecer también en Referencias.
	\end{alertblock}
\end{frame}

% ==================== SECCIÓN 7: APÉNDICES ====================
\section{Apéndices}

\begin{frame}{Apéndices}
	\begin{block}{Características}
		\begin{itemize}
			\item Título: ``Apéndices'' (si hay varios) o ``Apéndice'' (si hay uno solo)
			\item Nomenclatura: Apéndice A, Apéndice B, Apéndice C, etc.
			\item Cada apéndice inicia en página nueva
			\item Título centrado, negritas, mayúsculas y minúsculas
			\item Contenido justificado a la izquierda
		\end{itemize}
	\end{block}
	
	\begin{exampleblock}{Ejemplo de Título de Apéndice}
		\centering
		\textbf{Apéndice A}
		
		\textbf{Declaración de Inteligencia Artificial}
	\end{exampleblock}
	
	\begin{alertblock}{Tablas y Figuras en Apéndices}
		Numeración con letra del apéndice: Tabla A1, Figura B2, etc.
	\end{alertblock}
\end{frame}

\begin{frame}{Contenido de los Apéndices}
	\begin{block}{¿Qué Incluir?}
		\begin{itemize}
			\item Material complementario extenso
			\item Instrumentos de medición completos
			\item Listados detallados
			\item Datos crudos o tabulaciones extensas
			\item Documentos de respaldo
			\item Declaración de uso de IA (Apéndice A recomendado)
		\end{itemize}
	\end{block}
	
	\begin{alertblock}{Regla Importante}
		\begin{itemize}
			\item Cada apéndice debe tener su llamado en el texto
			\item Ejemplo: ``(ver Apéndice C)''
			\item Se presentan en el orden mencionado en el documento
		\end{itemize}
	\end{alertblock}
\end{frame}

% ==================== SECCIÓN 8: RECOMENDACIONES FINALES ====================
\section{Recomendaciones Finales}

\begin{frame}{Calidad de la Redacción}
	\begin{block}{Aspectos Fundamentales}
		\begin{itemize}
			\item \textbf{Lenguaje formal:} No coloquial, sin jergas ni eufemismos
			\item \textbf{Persona:} Evitar primera persona (singular y plural)
			\item \textbf{Párrafos:} Mínimo 4 oraciones (4 líneas)
			\item \textbf{Estructura:} Oración = mayúscula inicial + punto final
			\item \textbf{Puntuación:} Un espacio después de punto u otra puntuación
			\item \textbf{Coherencia:} Conectar ideas lógicamente
		\end{itemize}
	\end{block}
	
	\begin{alertblock}{Nunca}
		\begin{itemize}
			\item Subrayar texto
			\item Usar negritas (excepto títulos/subtítulos y tablas/figuras)
			\item Dejar espacios en blanco innecesarios
			\item Presentar documentos sin revisar ortografía y gramática
		\end{itemize}
	\end{alertblock}
\end{frame}

\begin{frame}{Estructura de Capítulos y Conclusiones}
	\begin{block}{Extensión Mínima}
		\begin{itemize}
			\item Trabajos aplicativos finales, trabajos de investigación final y tesis: mínimo 4 páginas por capítulo
			\item Cada capítulo inicia en página nueva
			\item Cada capítulo tiene conclusiones al final
		\end{itemize}
	\end{block}
	
	\begin{block}{Conclusiones de Capítulo}
		Las conclusiones son:
		\begin{itemize}
			\item Resultado del análisis profundo de la información
			\item Con sustentos académicos
			\item NO un resumen del capítulo
			\item Deben mostrar pensamiento crítico
		\end{itemize}
	\end{block}
\end{frame}

\begin{frame}{Plagio y Ética Académica}
	\begin{alertblock}{Definiciones Críticas}
		\begin{itemize}
			\item \textbf{Plagio:} Usar ideas/palabras ajenas sin dar crédito
			\item \textbf{Autoplagio:} No citar trabajos propios previos
			\item \textbf{Plagio involuntario:} Citar incorrectamente o no citar
		\end{itemize}
	\end{alertblock}
	
	\begin{block}{Cómo Evitar el Plagio}
		\begin{enumerate}
			\item Citar \alert{todas} las fuentes consultadas
			\item Parafrasear correctamente (no solo cambiar palabras)
			\item Usar comillas en citas textuales
			\item Incluir número de página en citas textuales
			\item Verificar que toda cita tenga su referencia
			\item Usar fuentes académicas confiables
		\end{enumerate}
	\end{block}
\end{frame}

\begin{frame}{Fuentes Académicas Confiables}
	\begin{block}{Bases de Datos Recomendadas}
		\begin{itemize}
			\item \textbf{Scopus} (multidisciplinaria)
			\item \textbf{Web of Science} (multidisciplinaria)
			\item \textbf{EBSCO} (negocios, salud, humanidades)
			\item \textbf{ProQuest} (tesis y disertaciones)
			\item \textbf{JSTOR} (humanidades y ciencias sociales)
			\item Libros académicos de editoriales reconocidas
			\item Artículos de revistas indexadas
		\end{itemize}
	\end{block}
	
	\begin{alertblock}{Evitar}
		\begin{itemize}
			\item Wikipedia como fuente primaria
			\item Blogs no académicos
			\item Fuentes sin autor identificable
			\item Sitios web sin respaldo institucional
		\end{itemize}
	\end{alertblock}
\end{frame}

\begin{frame}{Pies de Página}
	\begin{block}{Uso Limitado}
		Los pies de página en estilo APA se usan \alert{excepcionalmente}:
		\begin{itemize}
			\item Para complementar información que no puede presentarse de otra forma
			\item Para aclaraciones breves
			\item NO para referencias bibliográficas
		\end{itemize}
	\end{block}
	
	\begin{exampleblock}{Formato}
		\begin{itemize}
			\item Interlineado sencillo
			\item Fuente más pequeña (según configuración predeterminada)
			\item Numeración automática
		\end{itemize}
	\end{exampleblock}
	
	\begin{alertblock}{Recomendación}
		Minimizar su uso. Preferir integrar la información en el texto principal.
	\end{alertblock}
\end{frame}

\begin{frame}{Lista de Verificación Final}
	\begin{block}{Antes de Entregar, Verificar:}
		\small
		\begin{enumerate}
			\item ¿Portada completa y correcta?
			\item ¿Declaración Jurada de Autenticidad incluida?
			\item ¿Resumen/Abstract de 150-250 palabras?
			\item ¿Tabla de Contenidos actualizada?
			\item ¿Lista de Tablas y Figuras completas?
			\item ¿Capítulos con formato correcto (encabezados, seriación)?
			\item ¿Todas las citas con sus referencias?
			\item ¿Todas las referencias citadas en el texto?
			\item ¿Tablas y figuras numeradas y citadas correctamente?
			\item ¿Apéndices en orden y con llamados?
			\item ¿Revisión de ortografía y gramática completa?
			\item ¿Formato consistente en todo el documento?
		\end{enumerate}
	\end{block}
\end{frame}

% ==================== RECURSOS Y REFERENCIAS ====================
\section{Recursos Adicionales}

\begin{frame}{Referencias de Esta Presentación}
	\footnotesize
	\begin{itemize}
		\item American Psychological Association [APA]. (2020). \textit{Publication manual of the American Psychological Association} (7th ed.). https://doi.org/10.1037/0000165-000
		
		\item American Psychological Association [APA]. (2025). \textit{APA Style Blog}. https://apastyle.apa.org/blog
		
		\item D'Alessio, F. (2016). \textit{Observaciones generales para la redacción de trabajos académicos} (7a ed.). CENTRUM Católica Graduate Business School.
		
		\item Rojas, K., \& Marquina, P. S. (2021). \textit{Manual Básico de Estilo}. Centrum PUCP.
		
		\item Rojas Valdez, K. (2025). \textit{Manual para el Tesista: Guía del Estilo APA}. Centrum PUCP. http://xxx
	\end{itemize}
\end{frame}

\begin{frame}{Recursos en Línea}
	\begin{block}{Sitios Web Oficiales}
		\begin{itemize}
			\item \textbf{APA Style:} https://apastyle.apa.org/
			\item \textbf{APA Style Blog:} https://apastyle.apa.org/blog
			\item \textbf{Purdue OWL:} https://owl.purdue.edu/owl/research\_and\_citation/apa\_style/
		\end{itemize}
	\end{block}
	
	\begin{block}{Herramientas Útiles}
		\begin{itemize}
			\item Gestores bibliográficos: Mendeley, Zotero, EndNote
			\item Verificadores de plagio: Turnitin (institucional)
			\item Correctores ortográficos y gramaticales
		\end{itemize}
	\end{block}
\end{frame}

\begin{frame}{Preguntas Frecuentes - Parte 1}
	\begin{block}{¿Cuándo uso cursiva vs. comillas?}
		\begin{itemize}
			\item \textbf{Cursiva:} Títulos de libros, revistas, películas, tablas/figuras
			\item \textbf{Comillas:} Títulos de artículos, capítulos, citas textuales cortas
		\end{itemize}
	\end{block}
	
	\begin{block}{¿Qué hago si no encuentro el DOI?}
		Usar la URL del artículo. Preferir DOI cuando esté disponible.
	\end{block}
	
	\begin{block}{¿Incluyo bibliografía no citada?}
		NO. Solo incluir en Referencias lo que se citó en el texto.
	\end{block}
\end{frame}

\begin{frame}{Preguntas Frecuentes - Parte 2}
	\begin{block}{¿Cómo cito una comunicación personal?}
		En el texto: (N. Apellido, comunicación personal, fecha exacta)
		
		NO se incluye en Referencias.
	\end{block}
	
	\begin{block}{¿Puedo usar ``nosotros'' en mi tesis?}
		NO. Evitar primera persona. Usar voz pasiva o construcciones impersonales.
	\end{block}
	
	\begin{block}{¿Cuántas citas son suficientes?}
		Depende del contenido. Citar siempre que uses ideas ajenas. Evitar excesos y carencias.
	\end{block}
\end{frame}

% ==================== CIERRE ====================
\begin{frame}[plain,standout]
	\begin{center}
		{\Huge Gracias por su atención}
		
		\vspace{1cm}
		
		{\Large Preguntas y Comentarios}
		
		\vspace{1.5cm}
		
		\textbf{Centrum PUCP}\\
		Centro de Negocios\\
		Pontificia Universidad Católica del Perú
		
		\vspace{0.5cm}
		
		\small
		Para consultas adicionales sobre el estilo APA,\\
		revisar el \textit{Manual para el Tesista: Guía del Estilo APA}\\
		disponible en la biblioteca digital de Centrum PUCP
	\end{center}
\end{frame}

\end{document}